\documentclass[10pt,a4paper]{amsart}
\usepackage[margin=19mm]{geometry}
\usepackage{amsmath,amssymb,mathtools}
\usepackage[scr=rsfs,cal=euler]{mathalfa}
\usepackage{xfrac}
\usepackage[unicode,hidelinks,bookmarks=false]{hyperref}
\newcommand{\ie}{i.e.\ }
\newcommand{\cf}{cf.\ }
\newcommand{\resp}{respectively\ }
\newcommand{\Subsection}{\subsection}
\providecommand{\nobreakdash}{\nobreak\hbox{-}}
\setlength{\emergencystretch}{2em}
\multlinegap0pt
\usepackage{bm}
\usepackage{amssymb}
\usepackage{mathtools}
\usepackage{float}
\usepackage{xcolor}
\usepackage{tikz}
\usepackage{algorithm}\usepackage{algpseudocode}
\newtheorem{definition}{Definition}
\newtheorem{theorem}{Theorem}
\newcommand{\cD}{\mathcal{D}}
\newcommand{\cK}{\mathcal{K}}
\newcommand{\cN}{\mathcal{N}}
\newcommand{\cT}{\mathcal{T}}
\newcommand{\e}{\mathrm{e}} % Euler's constant
\newcommand{\inner}[2]{\langle#1, #2\rangle}
\title{An Optimal Control Approach to Transformer Training}
\author{Ka\u{g}an Akman, Naci Sald{\i}, Serdar Y\"uksel}
\date{Source: arXiv:2603.09571v1 (CC BY 4.0)}
\begin{document}
\maketitle

\noindent\emph{Source and licence.} An Optimal Control Approach to Transformer Training. Version arXiv:2603.09571v1 (CC BY 4.0), distributed by arXiv under the Creative Commons Attribution 4.0 International licence, http://creativecommons.org/licenses/by/4.0/, as recorded in the arXiv licence metadata for this exact version and shown on the abstract page https://arxiv.org/abs/2603.09571v1. Full licence notice: \texttt{licenses/arxiv-2603.09571-CC-BY-4.0.txt}.\par\smallskip

\noindent\textit{Editorial scope.} Counted unit: the article's theorem theorem (source0.tex lines 593--595). The article's complete original proof of it is delivered with it (source0.tex lines 596--609, one \texttt{proof} environment of 14 lines). No mathematics is written, repaired or re-derived: the statement and the proof are the article's own lines, in the article's own notation, and no retained line is substituted or re-flowed.  Retained uncounted as the source-local context the proof needs: lines 321--326; lines 391--395; lines 424--426.  Nothing was omitted from any retained span, so the package carries no omission record.  No retained line contains a citation, so the wrapper carries no bibliography and no bibliography entry.  No retained line contains a figure environment, an image-inclusion command or drawing code, so no image is delivered and the package declares no asset dependency.  Editorial formatting only, in addition: an AMS \texttt{amsart} preamble that restates the article's own declarations and macros used by the retained lines, namely the environments \texttt{definition}, \texttt{theorem} on separate counters, together with the standard packages those lines need.  The retained environments print in this document as Definition 1, Theorem 1 in order of appearance, whereas the article prints the same items with the numbering of its own full document.  The complete native source of the article as distributed in the arXiv e-print for this exact version is bundled unchanged as \texttt{sources/196072/source0.tex}.
\begin{equation}\label{eq:dynamics}
    \begin{cases}
        \displaystyle x_{t+1}^i = W_t\sigma(A_tx_t^i + b_t) + \sum_{j=1}^N \frac{ \e^{\beta\inner{Q_tx_t^i}{K_tx_t^j}}}{\sum_{j^\prime=1}^N\e^{\beta\inner{Q_tx_t^i}{K_tx_t^{j^\prime}}}}V_tx_t^j & \text{where } t\in\cT \\
        \displaystyle x^i_0 = x^i \in \mathrm{S}.
    \end{cases}
\end{equation}

\begin{align*}
    &\underset{\displaystyle\underline{\gamma}\in\Gamma}{\text{minimize}}\quad\frac{1}{NK}\sum_{i=1}^N\sum_{k=1}^K L(x_{T,\underline{\gamma}}^{i,k}, y^{i,k}) \\
    &\text{subject to: } x_{t,\underline{\gamma}}^{i,k} \text{ evolving according to \eqref{eq:dynamics}},\\
    &\hspace{5em} x_{0,\underline{\gamma}}^{i,k} = x^{i,k} \text{ for all } i\in\cN, k\in\cK, t\in\cT. \tag{OP}\label{prob:original}
\end{align*}

\begin{definition}\label{def:optimality}
    A strategy $\underline{\gamma}^*\in\Gamma$ is called \emph{globally optimal} if it attains the minimum in (\ref{prob:original}).
\end{definition}

\begin{theorem}
    If $\lambda > \frac{N}{2}\operatorname{diam}(\mathrm{S}, \|\cdot\|_2)^2$, then the lifted problem and (OP) are equivalent.
\end{theorem}

\begin{proof}
    Let $\bm{\mu}_{T,\underline{\gamma}} = (\mu_{T,\underline{\gamma}}^1,\ldots,\mu_{T,\underline{\gamma}}^K)$ be a generic terminal state evolved with $\bm{\Phi}$ under strategy $\underline{\gamma}\in\Gamma$ and initialized from the data set $\cD$. Respectively, let $\bm{\nu} := (\nu^1,\ldots,\nu^K)$ be the lifted labels, that is $\nu^k := \frac{1}{N}\sum_{i=1}^N\delta_{p_i, y^{i,k}}$ for all $k\in\cK$.

    In our case, where we compare empirical measures with the same number of atoms, the Wasserstein distance reduces to finding an optimal permutation. In particular, the optimal cost functional takes the form
    \begin{align*}
        V^*(\cD) &:= \inf_{\underline{\gamma}\in\Gamma}\frac{1}{K}\sum_{k=1}^K W_{2,\lambda}(\mu_{T,\underline{\gamma}}^k,\nu^k)\\
        &=\inf_{\underline{\gamma}\in\Gamma}\frac{1}{K}\sum_{k=1}^K \inf_{\sigma\in S_N}\frac{1}{N}\sum_{i=1}^N(\|x_{T,\underline{\gamma}}^{i,k} - y^{\sigma(i), k}\|_2^2 + \lambda|p_i - p_{\sigma(i)}|^2)
    \end{align*}
    Due to the choice of $\lambda$, we must have $\sigma$ as the identity permutation. Thus, we get
    \begin{equation*}
        V^*(\cD) = \inf_{\underline{\gamma}\in\Gamma}\frac{1}{NK}\sum_{k=1}^K \|x_{T,\underline{\gamma}}^{i,k} - y^{i, k}\|_2^2.
    \end{equation*}
    Recalling Definition \ref{def:optimality} concludes the proof.
\end{proof}

\end{document}
